\section{Discussion and Limitations}
\label{sec:discussion}

\textbf{Scope of field-centric rules.}
\tool{} is designed for requirements that can be grounded in protocol fields,
their conditions, and their implementation behavior. This scope includes many
parser, validation, negotiation, and error-handling rules. It does not replace
cryptographic proofs, full state-machine verification, or temporal reasoning
over unbounded sessions. Some cross-session properties may be approximated by
field rules, but they require a separate formal treatment before strong claims
can be made.

\textbf{Natural-language ambiguity.}
Standards contain optional behavior, informative examples, and modal verbs with
different force. \tool{} records modality and provenance, but final judgment may
still require expert interpretation for \emph{SHOULD}, \emph{MAY}, or
implementation-defined behavior. Our artifact therefore exposes the source span
and rationale for each rule rather than reporting a bare bug label.

\textbf{LLM dependence.}
\tool{} uses LLMs for extraction, ranking, semantic filtering, reasoning, and
test planning. The framework constrains these calls with schemas, field spaces,
syntax evidence, and deterministic validators, but it does not eliminate model
error. We mitigate this risk by retaining uncertainty, using low-temperature
decoding, validating outputs, comparing multiple models in the evaluation, and
requiring executable evidence for confirmed reports.

\textbf{Closed-source and generated implementations.}
The prototype assumes source access and at least partial parser support. Closed
implementations, heavily generated code, and macro-heavy systems reduce the
quality of symbol indexing and call-graph recovery. Black-box testing can
partially compensate, but specification-to-variable alignment becomes weaker.

\textbf{Ground truth.}
Protocol conformance is not always a binary property. Some standards are
ambiguous, some implementations intentionally deviate for compatibility, and
some rules describe recommendations rather than obligations. We therefore use a
state-based adjudication protocol: reports are submitted, confirmed/fixed,
rejected as false positives, or kept pending until maintainer feedback,
independent review, or reproducible tests provide enough evidence. This protocol
supports adjudicated precision but not a whole-corpus recall claim.

% \textbf{Cost and scalability.}
% Multi-round reasoning and validation are more expensive than one-shot
% prompting. \tool{} reduces cost by caching standard extraction, limiting code
% contexts before LLM calls, and invoking validation only for suspicious or
% unresolved samples. Still, deployment on very large protocol ecosystems
% may require prioritization. One practical strategy is to rank rules by modality,
% security relevance, and historical bug density, then validate high-risk rules
% first.

\textbf{Threats to validity.}
Our evaluation has three main threats. First, benchmark selection may
overrepresent well-structured open-source implementations and underrepresent
closed or generated systems. Second, manually labeled ground truth may reflect
reviewer interpretation of ambiguous standards. Third, LLM nondeterminism can
change extraction or reasoning outcomes across runs. We mitigate these threats
by reporting subject-selection criteria, using two-author adjudication,
preserving raw evidence, separating pending reports from confirmed fixes, and
measuring variance across repeated GPT-5.4 runs and model families.

% \textbf{What would make a report actionable?}
% Maintainers are unlikely to act on a vague statement that code and standard
% ``seem inconsistent.'' A useful report should include the exact standard span,
% the normalized rule, the relevant field and implementation variable, the code
% path that accepts or rejects the input, the build configuration, and a minimal
% reproduction script. \tool{} emits this bundle and stores it under a content
% hash so that reviewers can inspect the evidence without rerunning the full
% pipeline.

\textbf{Generalization beyond protocols.}
The field-centric idea may transfer to other domains where natural-language
documents define structured objects, such as file formats, API contracts, or
hardware specifications. The representation would need different object spaces:
API parameters instead of protocol fields, register bits instead of message
fields, or configuration keys instead of transport parameters. This paper
focuses on network protocols because their standards and implementations make
the alignment problem both important and concrete.

\textbf{Ethical considerations.}
We disclosed confirmed security-relevant findings to maintainers before
publication. The public artifact avoids publishing exploit-ready inputs for
unpatched vulnerabilities while still providing enough artifact detail for
reviewers to validate the claims under confidentiality constraints.
